\documentclass[12pt,a4paper]{article}

% --- Pakiety dla języka polskiego ---
\usepackage[utf8]{inputenc}
\usepackage[T1]{fontenc}
\usepackage[polish]{babel}
\usepackage{tabularx}
\usepackage{booktabs}
\usepackage{listings}
\usepackage{xcolor}
\usepackage[margin=2.5cm]{geometry}
\usepackage{enumitem}
\usepackage{hyperref}
\hypersetup{
    colorlinks=false,       % false: ramki; true: kolorowy tekst
    pdfborder={0 0 0},     % ustawia szerokość ramki na 0 (całkowicie znika)
}

% --- Konfiguracja pakietu listings dla polskich znaków ---
\title{Dokumentacja Implementacyjna Projektu: \\ Generowanie Grafów Planarnych}
\author{Aleksandra Rybałko, Nella Karczewicz}
\date{\today}

\begin{document}

\maketitle

% --- DODANY SPIS TREŚCI ---
\tableofcontents
\newpage
% --------------------------

% --- Konfiguracja pakietu listings dla polskich znaków ---
\lstset{
    backgroundcolor=\color{gray!5},
    basicstyle=\ttfamily\small,
    breaklines=true,
    captionpos=b,
    commentstyle=\color{green!60!black},
    keywordstyle=\color{blue},
    stringstyle=\color{orange},
    frame=single,
    showstringspaces=false,
    % Rozwiązanie problemu polskich znaków:
    literate={ą}{{\k{a}}}1 {ć}{{\'c}}1 {ę}{{\k{e}}}1 {ł}{{\l{}}}1 {ń}{{\'n}}1 {ó}{{\'o}}1 {ś}{{\'s}}1 {ź}{{\'z}}1 {ż}{{\.z}}1
             {Ą}{{\k{A}}}1 {Ć}{{\'C}}1 {Ę}{{\k{E}}}1 {Ł}{{\L{}}}1 {Ń}{{\'N}}1 {Ó}{{\'O}}1 {Ś}{{\'S}}1 {Ź}{{\'Z}}1 {Ż}{{\.Z}}1
}



\section{Cel i zakres implementacji}
Celem niniejszej dokumentacji jest przedstawienie szczegółów technicznych oraz architektury programu napisanego w języku C, służącego do automatycznego rozmieszczania wierzchołków grafów planarnych. Implementacja koncentruje się na zapewnieniu stabilności obliczeń fizycznych, efektywnym zarządzaniu pamięcią oraz rygorystycznej walidacji danych wejściowych.

\section{Architektura systemu i struktury danych}

Aplikacja została zaprojektowana w oparciu o paradygmat programowania modularnego. Pozwala to na ścisłą separację logiki wczytywania danych, zarządzania strukturami oraz algorytmów rozmieszczania wierzchołków.

\subsection{Podział na moduły funkcjonalne}
Kod źródłowy został podzielony na niezależne jednostki logiczne, co ułatwia konserwację i niezależne testowanie komponentów:

\begin{itemize}[label=\textbullet]
    \item \textbf{\texttt{graph.h}:} Centralny moduł definicji struktur. Zawiera opisy typów \texttt{Node}, \texttt{Edge} oraz \texttt{Graph}, stanowiące wspólny interfejs dla całego systemu.
    \item \textbf{\texttt{io.c / io.h}:} Moduł wejścia/wyjścia. Odpowiada za niskopoziomową obsługę plików, dwufazową walidację składniową, sanityzację danych oraz eksport wyników (TXT/BIN).
    \item \textbf{\texttt{layout\_tutte.c / layout\_tutte.h}:} Implementacja deterministycznego algorytmu barycentrycznego oraz geometrycznych testów przecięć krawędzi (CCW).
    \item \textbf{\texttt{fruchterman.c / fruchterman.h}:} Implementacja heurystycznego modelu siłowego wraz z mechanizmami grawitacji i chłodzenia układu.
    \item \textbf{\texttt{main.c}:} Moduł sterujący. Zarządza logiką CLI (\textit{Command Line Interface}), koordynuje przepływ danych i realizuje pętlę wielokrotnych prób dla algorytmów probabilistycznych.
\end{itemize}

\subsection{Reprezentacja grafu w pamięci operacyjnej}
Program wykorzystuje trzy kluczowe struktury zdefiniowane w pliku \texttt{graph.h}.

\subsubsection{Wierzchołek (\texttt{Node})}
Struktura ta przechowuje pełny stan fizyczny pojedynczego punktu:
\begin{lstlisting}[language=C, caption={Struktura wierzchołka w graph.h}]
typedef struct {
    int id;          // Oryginalny identyfikator z pliku
    double x, y;     // Aktualne współrzędne położenia
    double dx, dy;   // Wektory sił (akumulatory przesunięcia)
} Node;
\end{lstlisting}
\textbf{Analiza:} Zastosowanie pól \texttt{dx} i \texttt{dy} wewnątrz struktury pozwala na uniknięcie dodatkowych alokacji tablic pomocniczych podczas iteracji algorytmu Fruchtermanna, co znacząco poprawia wydajność procesora.

\subsubsection{Krawędź (\texttt{Edge})}
Krawędzie definiują relacje sąsiedztwa w grafie:
\begin{lstlisting}[language=C, caption={Struktura krawędzi w graph.h}]
typedef struct {
    int u_idx;       // Indeks wierzchołka A w głównej tablicy
    int v_idx;       // Indeks wierzchołka B w głównej tablicy
    double weight;   // Waga połączenia
} Edge;
\end{lstlisting}
\textbf{Analiza:} Przechowywanie indeksów (\texttt{u\_idx}, \texttt{v\_idx}) zamiast surowych identyfikatorów pozwala na dostęp do danych sąsiadów w czasie stałym $O(1)$, co jest niezbędne dla szybkości działania metody barycentrycznej.

\subsubsection{Główny kontener (\texttt{Graph})}
Agreguje wszystkie zasoby grafu w jednym obiekcie:
\begin{lstlisting}[language=C, caption={Struktura grafu w graph.h}]
typedef struct {
    Node *nodes;      // Dynamiczna tablica wierzchołków
    Edge *edges;      // Dynamiczna tablica krawędzi
    int node_count;   // Liczba wczytanych wierzchołków
    int edge_count;   // Liczba wczytanych krawędzi
    double width;     // Szerokość obszaru (domyślnie 1000.0)
    double height;    // Wysokość obszaru (domyślnie 1000.0)
} Graph;
\end{lstlisting}
\newpage
\subsection{Schemat przepływu danych i cykl życia programu}
Proces przetwarzania danych jest zorganizowany w sposób kaskadowy, z wbudowanymi punktami kontrolnymi (tzw. "strażnikami"), które mogą przerwać wykonanie w przypadku wykrycia błędów krytycznych.

\begin{enumerate}
    \item \textbf{Inicjalizacja CLI:} Moduł \texttt{main.c} przetwarza argumenty użytkownika, określając ścieżki plików oraz wybór algorytmu.
    \item \textbf{Walidacja Diagnostyczna:} Moduł \texttt{io.c} skanuje plik wejściowy bez alokacji struktur, raportując ewentualne błędy formatu w konkretnych liniach.
    \item \textbf{Budowa Topologii:} Po sukcesie walidacji następuje dynamiczna alokacja pamięci w strukturze \texttt{Graph} i mapowanie identyfikatorów zewnętrznych na indeksy tablicy.
    \item \textbf{Weryfikacja Topologiczna:} Program sprawdza spójność grafu (DFS) oraz potencjalną planarność (wzór Eulera).
    \item \textbf{Obliczenia Geometryczne:} Wybrany moduł (\texttt{Tutte} lub \texttt{Fruchterman}) modyfikuje współrzędne \texttt{x, y} w strukturach \texttt{Node}.
    \item \textbf{Eksport i Sprzątanie:} Wyniki są zapisywane do pliku, a pamięć dynamiczna jest zwalniana przez funkcję \texttt{free\_graph}.
\end{enumerate}

\section{Zarządzanie pamięcią i cykl życia obiektu}
Program realizuje model bezpiecznego zarządzania pamięcią, dostosowując rozmiar struktur do rzeczywistej wielkości grafu.

\subsection{Alokacja dynamiczna}
Pamięć jest rezerwowana wewnątrz modułu \texttt{io.c} po zakończeniu fazy walidacji pliku:
\begin{lstlisting}[language=C, caption={Alokacja w funkcji read\_graph}]
g->edges = malloc(1000 * sizeof(Edge)); 
g->nodes = calloc(1000, sizeof(Node));
\end{lstlisting}
\textbf{Analiza:} Wykorzystanie \texttt{calloc} gwarantuje domyślną inicjalizację pól zerem, co zapobiega błędom "garbage values" przy startowaniu symulacji fizycznej.

\subsection{Zwalnianie zasobów}
Za czyszczenie pamięci odpowiada funkcja \texttt{free\_graph}, wywoływana w każdym scenariuszu zakończenia pracy programu:
\begin{lstlisting}[language=C, caption={Funkcja zwalniająca miejsce w pamięci}]
void free_graph(Graph *g) {
    if(g->nodes) free(g->nodes);
    if(g->edges) free(g->edges);
    g->nodes = NULL; 
    g->edges = NULL;
}
\end{lstlisting}
\section{Moduł wejścia/wyjścia i rygorystyczna walidacja danych}

Moduł \texttt{io.c} stanowi fundament stabilności aplikacji. Jego zadaniem jest nie tylko odczyt danych, ale przede wszystkim ich filtracja, ujednolicenie oraz weryfikacja założeń topologicznych grafu przed etapem obliczeniowym.

\subsection{Automatyczna detekcja formatu i separatorów}
Program został wyposażony w inteligentny mechanizm rozpoznawania struktury pliku, co pozwala na obsługę różnych standardów zapisu danych (CSV, SSV).

\begin{lstlisting}[language=C, caption={Funkcja detect\_separator}]
static char detect_separator(const char *line) {
    return (strchr(line, ';') == NULL && strchr(line, ',') != NULL) ? ',' : ';';
}
\end{lstlisting}
\textbf{Wyjaśnienie logiki:} Funkcja analizuje surowy ciąg znaków w poszukiwaniu separatorów. Priorytetyzacja średnika nad przecinkiem zapobiega błędnej interpretacji w przypadku, gdy przecinek pełni rolę separatora dziesiętnego w wagach krawędzi.

\subsection{Sanityzacja danych i obsługa błędów lokalizacji}
Ze względu na różnice w zapisie liczb zmiennoprzecinkowych (kropka vs przecinek), zaimplementowano funkcję czyszczącą.

\begin{lstlisting}[language=C, caption={Funkcja sanitize w io.c}]
int sanitize(char *s) {
    int changed = 0;
    for (int i = 0; s[i]; i++) {
        if (s[i] == ',') {
            s[i] = '.';
            changed = 1;
        }
    }
    return changed;
}
\end{lstlisting}
\textbf{Wyjaśnienie logiki:} Funkcja iteruje po napisie i zamienia każdy przecinek na kropkę. Jest to kluczowe dla poprawnego działania funkcji \texttt{sscanf}, która w standardzie C oczekuje kropki jako separatora dziesiętnego.
\newpage
\subsection{Zapis danych w formacie binarnym i tekstowym}
Funkcja \texttt{save\_graph} pozwala na eksport współrzędnych w zależności od potrzeb użytkownika.
\begin{lstlisting}[language=C, caption={Obsługa formatu binarnego w save\_graph}]
if (strcmp(format, "bin") == 0) {
    fwrite(&(g->node_count), sizeof(int), 1, f);
    
    fwrite(g->nodes, sizeof(Node), g->node_count, f);
    }
    \end{lstlisting}
\textbf{Wyjaśnienie logiki:} W trybie binarnym dane są zapisywane bezpośrednio z pamięci RAM przy użyciu \texttt{fwrite}. Pozwala to na uniknięcie błędów zaokrągleń oraz drastycznie zmniejsza rozmiar pliku wyjściowego w porównaniu do zapisu tekstowego \texttt{fprintf}.


\section{Założenia teoretyczne i walidacja topologiczna}

Aby algorytmy rozmieszczania wierzchołków mogły wygenerować poprawny i czytelny rysunek, graf wejściowy musi spełniać określone założenia matematyczne. Program weryfikuje te warunki na etapie walidacji wstępnej, niezależnie od wybranego algorytmu.

\subsection{Fundamenty teoretyczne: Struktura grafu, Spójność i Planarność}
W celu zachowania stabilności obliczeń numerycznych oraz poprawności topologicznej, system operuje na tzw. \textbf{grafach prostych} i sprawdza ich kluczowe parametry:

\begin{itemize}
    \item \textbf{Wymóg grafu prostego:} Program automatycznie odrzuca \textbf{pętle własne} (połączenie wierzchołka z samym sobą) oraz ignoruje \textbf{krawędzie wielokrotne} (duplikaty między tymi samymi węzłami). Pętle własne uniemożliwiłyby działanie modelu siłowego ze względu na zerową odległość (dzielenie przez zero), a duplikaty zaburzałyby statyczną weryfikację planarności.
    \item \textbf{Spójność (Connectivity):} Jest to warunek konieczny dla algorytmów siłowych i barycentrycznych. W grafie spójnym istnieje ścieżka między każdą parą wierzchołków. Brak spójności spowodowałby "nakładanie się" osobnych komponentów na siebie lub ich niekontrolowane rozbieganie się w obszarze roboczym pod wpływem sił odpychania.
    \item \textbf{Planarność i wzór Eulera:} Graf jest planarny, jeśli można go narysować bez przecięć krawędzi. Wykorzystujemy tu wniosek z twierdzenia Eulera: dla grafu planarnego o $V \ge 3$ wierzchołkach, liczba krawędzi $E$ musi spełniać warunek $E \le 3V - 6$. Jest to pierwszy "strażnik" blokujący przetwarzanie zbyt gęstych danych.
\end{itemize}
\newpage
\subsection{Weryfikacja spójności grafu (Algorytm DFS)}
Zaimplementowano funkcję wykorzystującą przeszukiwanie w głąb (DFS) ze stosem, która zlicza odwiedzone węzły startując od pierwszego wczytanego punktu.

\begin{lstlisting}[language=C, caption={Funkcja is\_graph\_connected}]
int is_graph_connected(Graph *g) {
    if (g->node_count <= 1) return 1;
    int *visited = calloc(g->node_count, sizeof(int));
    int *stack = malloc(g->node_count * sizeof(int));
    int top = -1;

    stack[++top] = 0; // Start od pierwszego wierzchołka
    visited[0] = 1;
    int visited_count = 1;

    while (top >= 0) {
        int u = stack[top--];
        for (int i = 0; i < g->edge_count; i++) {
            int v = -1;
            if (g->edges[i].u_idx == u) v = g->edges[i].v_idx;
            else if (g->edges[i].v_idx == u) v = g->edges[i].u_idx;

            if (v != -1 && !visited[v]) {
                visited[v] = 1;
                stack[++top] = v;
                visited_count++;
            }
        }
    }
    free(visited); free(stack);
    return (visited_count == g->node_count);
}
\end{lstlisting}

\textbf{Logika blokady:} Jeśli $visited\_count < g->node\_count$, graf jest niespójny. Program przerywa pracę z kodem błędu 9, co zapobiega generowaniu błędnych wyników.

\subsection{Weryfikacja potencjalnej planarności (Wzór Eulera)}
Drugim etapem walidacji wspólnej jest statyczne sprawdzenie gęstości grafu przed uruchomieniem kosztownych obliczeń fizycznych.

\begin{lstlisting}[language=C, caption={Funkcja is\_potentially\_planar}]
int is_potentially_planar(Graph *g) {
    if (g->node_count <= 2) return 1;
    return (g->edge_count <= (3 * g->node_count - 6));
}
\end{lstlisting}

\textbf{Logika blokady:} Jeśli warunek nie zostanie spełniony, program informuje o braku potencjalnej planarności (kod błędu 7). Stanowi to filtr chroniący przed próbą siłowego rozmieszczania struktur, które matematycznie nie posiadają reprezentacji planarnej.
% --- Sekcja 5: Algorytm Tutte'a ---
\newpage
\section{Algorytm Tutte'go – Deterministyczna Metoda Barycentryczna}

Algorytm Tutte'go, znany również jako metoda barycentryczna, stanowi fundament deterministycznego podejścia do wizualizacji grafów planarnych. W przeciwieństwie do metod heurystycznych, algorytm ten opiera się na rozwiązaniu układu równań liniowych, co gwarantuje uzyskanie unikalnego i stabilnego wyniku dla danej topologii grafu.

\subsection{Logika i podstawy matematyczne}
Głównym założeniem algorytmu jest fakt, że każdy wierzchołek wewnętrzny powinien znajdować się w tzw. barycentrum (środku ciężkości) swoich sąsiadów. 

Zgodnie z twierdzeniem Tutte'a, aby uzyskać ułożenie planarne (brak przecięć krawędzi), muszą zostać spełnione trzy warunki:
\begin{enumerate}
    \item Graf musi być planarny i trójspójny.
    \item Wierzchołki zewnętrznej ściany grafu (rama) muszą tworzyć wielokąt wypukły.
    \item Każdy wierzchołek wewnętrzny musi być połączony z co najmniej trzema innymi wierzchołkami.
\end{enumerate}

Pozycja każdego ruchomego wierzchołka $v_i$ jest wyznaczana jako średnia arytmetyczna pozycji jego sąsiadów:
\[ x_i = \frac{1}{|Adj(v_i)|} \sum_{j \in Adj(v_i)} x_j, \quad y_i = \frac{1}{|Adj(v_i)|} \sum_{j \in Adj(v_i)} y_j \]

\subsection{Implementacja procesu inicjalizacji}
Pierwszym etapem algorytmu jest zdefiniowanie "kotwic" układu, czyli wierzchołków stałych, które nie podlegają procesowi relaksacji.

\begin{lstlisting}[language=C, caption={Inicjalizacja ramy grafu w layout\_tutte.c}]
void initialize_tutte_fixed_nodes(Graph *g) {
    if (g->node_count < 3) return;

    double centerX = g->width / 2.0;
    double centerY = g->height / 2.0;
    // Promień ramy dostosowany do wymiarów obszaru roboczego
    double radius = (g->width < g->height ? g->width : g->height) / 3.0;

    // Rozmieszczenie pierwszych 4 wierzchołków na obwodzie koła
    for (int i = 0; i < 4 && i < g->node_count; i++) {
        double angle = 2.0 * M_PI * i / 4.0;
        g->nodes[i].x = centerX + radius * cos(angle);
        g->nodes[i].y = centerY + radius * sin(angle);
    }
}
\end{lstlisting}
\textbf{Wyjaśnienie logiki:} Program rezerwuje indeksy $0-3$ jako wierzchołki zewnętrzne. Wykorzystanie funkcji \texttt{cos} i \texttt{sin} zapewnia idealnie symetryczne, wypukłe ułożenie startowe, co jest warunkiem koniecznym dla powodzenia algorytmu.

\subsection{Iteracyjne rozwiązywanie układu (Metoda relaksacji)}
Zamiast budować złożone macierze rzadkie, w projekcie zastosowano metodę relaksacji, która w sposób przybliżony dąży do stanu równowagi.

\begin{lstlisting}[language=C, caption={Logika obliczania środka ciężkości}]
for (int iter = 0; iter < 200; iter++) {
    // Aktualizujemy tylko wierzchołki ruchome (od indeksu 4 wzwyż)
    for (int i = 4; i < g->node_count; i++) {
        double sum_x = 0, sum_y = 0;
        int neighbors_count = 0;

        // Przeszukiwanie listy krawędzi w celu znalezienia sąsiadów wierzchołka i
        for (int j = 0; j < g->edge_count; j++) {
            int neighbor_idx = -1;
            if (g->edges[j].u_idx == i) neighbor_idx = g->edges[j].v_idx;
            else if (g->edges[j].v_idx == i) neighbor_idx = g->edges[j].u_idx;

            if (neighbor_idx != -1) {
                sum_x += g->nodes[neighbor_idx].x;
                sum_y += g->nodes[neighbor_idx].y;
                neighbors_count++;
            }
        }
        // Przesunięcie wierzchołka do barycentrum jego sąsiedztwa
        if (neighbors_count > 0) {
            g->nodes[i].x = sum_x / neighbors_count;
            g->nodes[i].y = sum_y / neighbors_count;
        }
    }
}
\end{lstlisting}
\textbf{Wyjaśnienie logiki:} Pętla zewnętrzna wykonuje 200 iteracji, co empirycznie dobrano jako wartość zapewniającą stabilizację wizualną grafu. Wewnątrz, dla każdego węzła ruchomego, program sumuje pozycje sąsiadów i wyciąga średnią. Jeśli wierzchołek nie posiada sąsiadów (\texttt{neighbors\_count == 0}), pozostaje w miejscu, co zapobiega błędom dzielenia przez zero.
\newpage
% --- Sekcja 6: Fruchterman-Reingold ---
% --- Sekcja 6: Fruchterman-Reingold ---
\section{Algorytm Fruchterman-Reingold – Heurystyka Siłowa}

Algorytm Fruchterman-Reingold opiera się na analogii fizycznej, w której wierzchołki traktowane są jak cząsteczki atomowe, a krawędzie jak sprężyny.

\subsection{Logika oddziaływań fizycznych}
Głównym celem algorytmu jest znalezienie konfiguracji o najniższej energii. Kluczowym parametrem jest stała optymalnej odległości $k$:
\[ k = \sqrt{\frac{Width \cdot Height}{|V|}} \]
.

\subsection{Specyfika implementacji i dążenie do planarności}
W klasycznym ujęciu algorytm ten jest heurystyką estetycznego rozmieszczania, która dąży do równomiernego rozłożenia węzłów, ale nie gwarantuje uzyskania układu planarnego. W niniejszym projekcie algorytm został jednak rozbudowany o mechanizmy wymuszające bezkolizyjność:

\begin{itemize}
    \item \textbf{Zastosowanie grawitacji:} Dodanie wektora siły dociągającej do środka zapobiega nadmiernemu rozproszeniu wierzchołków.
    \item \textbf{Wielokrotne losowe starty:} Zaimplementowano pętlę \textit{Retry Logic} (do 3000 prób), co pozwala na uniknięcie lokalnych minimów energetycznych.
    \item \textbf{Weryfikacja testem przecięć:} Program uznaje działanie za zakończone sukcesem dopiero po uzyskaniu statusu \texttt{is\_layout\_planar == 1}.
\end{itemize}

\subsection{Implementacja sił odpychania i grawitacji}
W przeciwieństwie do sił przyciągania, siły odpychania działają na każdą parę wierzchołków w grafie, co nadaje algorytmowi złożoność $O(V^2)$.

\begin{lstlisting}[language=C, caption={Obliczanie sumarycznych wektorów sił}]
// 1. ODPYCHANIE: Każdy wierzchołek odpycha się od każdego innego
for (int i = 0; i < graph->node_count; i++) {
    graph->nodes[i].dx = 0; // Resetowanie wektora przesunięcia
    graph->nodes[i].dy = 0;

    // Grawitacja: słaba siła dociągająca wierzchołek do środka ekranu
    double dx_center = centerX - graph->nodes[i].x;
    double dy_center = centerY - graph->nodes[i].y;
    graph->nodes[i].dx += dx_center * 0.05; 
    graph->nodes[i].dy += dy_center * 0.05;

    for (int j = 0; j < graph->node_count; j++) {
        if (i == j) continue;
        double vx = graph->nodes[i].x - graph->nodes[j].x;
        double vy = graph->nodes[i].y - graph->nodes[j].y;
        \end{lstlisting}
        \begin{lstlisting}
        double dist = sqrt(vx * vx + vy * vy);

        if (dist < 0.01) dist = 0.01; // Unikanie osobliwości przy dist=0

        double force = (k * k) / dist; // f_r(d) = k^2 / d
        graph->nodes[i].dx += (vx / dist) * force;
        graph->nodes[i].dy += (vy / dist) * force;
    }
}
\end{lstlisting}
\textbf{Wyjaśnienie logiki:} Wektory \texttt{dx} i \texttt{dy} pełnią rolę akumulatorów sił. Dodanie grawitacji (siła dociągająca do środka) zapobiega "rozlatywaniu się" niespójnych komponentów grafu poza obszar roboczy.

\subsection{Mechanizm chłodzenia (Simulated Annealing)}
Aby układ mógł "zastygnąć" w stabilnej pozycji, wprowadzono parametr temperatury $t$, który ogranicza maksymalne przemieszczenie wierzchołka w jednej iteracji.

\begin{lstlisting}[language=C, caption={Aktualizacja pozycji i chłodzenie}]
for (int i = 0; i < graph->node_count; i++) {
    double disp_dist = sqrt(graph->nodes[i].dx * graph->nodes[i].dx + 
                           graph->nodes[i].dy * graph->nodes[i].dy);

    if (disp_dist > 0) {
        // Ograniczenie przesunięcia aktualną temperaturą t
        double limited_dist = (disp_dist < t) ? disp_dist : t;
        graph->nodes[i].x += (graph->nodes[i].dx / disp_dist) * limited_dist;
        graph->nodes[i].y += (graph->nodes[i].dy / disp_dist) * limited_dist;
    }
    // (...) ograniczenie do ramki obszaru
}
t *= 0.95; // Stopniowe chłodzenie systemu o 5% w każdym kroku
\end{lstlisting}
\textbf{Wyjaśnienie logiki:} Temperatura początkowa jest wysoka, co pozwala wierzchołkom na duże skoki i wyjście z lokalnych minimów energetycznych. Z każdą iteracją temperatura maleje (\texttt{t *= 0.95}), co powoduje, że graf staje się coraz bardziej "sztywny" i stabilny.


\newpage

\section{Scenariusze testowe}

W celu potwierdzenia niezawodności aplikacji przeprowadzono serię testów systemowych. Dokumentacja zestawia fragmenty plików wejściowych z logami z terminala, co pozwala na precyzyjną weryfikację logiki walidacyjnej modułu \texttt{io.c}.

\subsection{Szczegółowa analiza przypadków testowych (Pliki .txt)\\}

\subsubsection{Brak planarności – zbyt duża gęstość (\texttt{err\_gestosc.txt})}
W tej części testów poddano weryfikacji skuteczność mechanizmów kontrolnych programu w sytuacji, gdy wprowadzony graf jest matematycznie nieplanarny. W wyniku mamy porównanie, w jaki sposób obie zaimplementowane metody Tutta oraz Fruchtermana – reagują na brak możliwości uzyskania bezkolizyjnego ułożenia.
\begin{lstlisting}[language=bash, caption={Log terminala - wynik działania algorytmu Tutte}]
Uruchamiam algorytm Tutte...

Sprawdzanie planarności wygenerowanego układu (Tutte)...

============================================================
   PORAŻKA: WYNIK TUTTE NIE JEST PLANARNY    
============================================================
Zapis pliku został zablokowany.


\end{lstlisting}
\begin{lstlisting}[language=bash, caption={Log terminala - wynik działania algorytmu Fruchtermana}]
Uruchamiam symulację Fruchtermana (max 3000 prób)...

Sprawdzanie planarności wygenerowanego układu (Fruchterman)...

============================================================
   KOMUNIKAT: NIE ZNALEZIONO UKŁADU PLANARNEGO    
============================================================
Po 3000 prób algorytm FR nie wyeliminował przecięć.
Zapis pliku wyjściowego anulowany. Spróbuj ponownie.


\end{lstlisting}
\newpage
\subsubsection{Błędny format danych (\texttt{err\_litera.txt})}
Test sprawdza odporność walidatora na znaki niebędące cyframi w polach numerycznych (U, V, Waga).

\textbf{Fragment struktury:}
\begin{verbatim}
krawedz1;1;2x;1.O
krawedz2;2;X;1.5
krawedz3;3;4;0.8
\end{verbatim}

\begin{lstlisting}[language=bash, caption={Log terminala - wykrycie liter}]
Uruchamiam przetwarzanie...

===============================================================
 RAPORT ANALIZY DANYCH WEJSCIOWYCH 
===============================================================
Linia    | Zinterpretowane dane                          | Status         
---------------------------------------------------------------

1        | [ !!! BŁĄD KRYTYCZNY - ZŁY FORMAT !!! ]    | BŁĄD         
2        | [ !!! BŁĄD KRYTYCZNY - ZŁY FORMAT !!! ]    | BŁĄD         
3        | Krawedz krawedz3: 3 -> 4 (w: 0.8)            | [OK]           
===============================================================

STOP! Znaleziono 2 bledow krytycznych w strukturze pliku.
 -> INFO: Poprawny format to: Nazwa;U;V;Waga
 -> Separator danych to ŚREDNIK (;), a separator dziesiętny to KROPKA (.)
 -> ID wierzchołków muszą mieścić się w przedziale od 1 do 1000.
 -> Waga musi być dodatnia, z jednym miejscem po kropce.

===============================================================
   PRZETWARZANIE PRZERWANE: BŁĘDNA STRUKTURA GRAFU    
===============================================================
Powód: Wykryto błędy w strukturze grafu (np. niepoprawny format 
lub niedozwolone pętle własne). Sprawdź raport powyżej i popraw
plik wejściowy.
===============================================================


\end{lstlisting}
\newpage
\subsubsection{Minimalna liczba wierzchołków (\texttt{err\_za\_maly.txt})}
Algorytm Tutte'a wymaga minimum 3 węzłów do stworzenia ramy. Test sprawdza blokadę dla zbyt małych struktur.

\textbf{Fragment struktury:}
\begin{verbatim}
k1;1;2;1.0
\end{verbatim}

\begin{lstlisting}[language=bash, caption={Log terminala - za mały graf}]
Uruchamiam przetwarzanie...

===============================================================
 RAPORT ANALIZY DANYCH WEJSCIOWYCH 
===============================================================
Linia    | Zinterpretowane dane                          | Status         
---------------------------------------------------------------
1        | Krawedz k1: 1 -> 2 (w: 1.0)                  | [OK]           
===============================================================

Wczytywanie zakończone sukcesem. Dane są poprawne technicznie.
Statystyki: 2 węzłów, 1 krawędzi.
---------------------------------------------------------------

============================================================
   BŁĄD: NIEWYSTARCZAJĄCA LICZBA ELEMENTÓW GRAFU
============================================================
Statystyki wczytanego pliku:
-> Wierzchołki: 2 (wymagane >= 3)
-> Krawędzie:   1 (wymagane >= 2)

POWÓD: Graf jest zbyt mały, aby wyznaczyć układ planarny.
Upewnij się, że plik [testy_N/err_za_maly.txt] zawiera poprawny graf.

============================================================

\end{lstlisting}
\newpage
\subsubsection{Niespójność grafu (\texttt{err\_niespojny.txt})}
Weryfikacja algorytmu DFS. Program musi wykryć, że graf składa się z kilku niepołączonych komponentów.

\textbf{Fragment struktury:}
\begin{verbatim}
k1;1;2;1.0
k2;2;3;1.0
k3;4;5;1.0
\end{verbatim}

\begin{lstlisting}[language=bash, caption={Log terminala - graf niespójny}]
Sprawdzanie spójności grafu za pomocą algorytmu DFS ...

============================================================
   BŁĄD KRYTYCZNY: GRAF NIESPÓJNY    
============================================================
Dane z pliku wejściowego [testy_N/err_niespojny.txt] są błędne:
-> Graf składa się z kilku niepołączonych ze sobą części.
-> Algorytm wymaga grafu spójnego, aby poprawnie wyznaczyć pozycje.
============================================================

\end{lstlisting}
\subsubsection{Autokorekta separatora dziesiętnego (\texttt{test\_autokorekta.txt})}
Weryfikacja funkcji \texttt{sanitize}, która zamienia przecinki na kropki, umożliwiając poprawny odczyt wag przez \texttt{sscanf}.

\textbf{Fragment struktury:} \begin{verbatim}
krawedz1;1;2;1,5
krawedz2;2;3;2,7
krawedz3;3;1;0,9
\end{verbatim}

\begin{lstlisting}[language=bash, caption={Log terminala - INFO: Dokonano autokorekt}]
Uruchamiam przetwarzanie...

===============================================================
 RAPORT ANALIZY DANYCH WEJSCIOWYCH 
===============================================================
Linia    | Zinterpretowane dane               | Status         
---------------------------------------------------------------
1        | Krawedz krawedz1: 1 -> 2 (w: 1.50) | [Poprawiono ,] 
2        | Krawedz krawedz2: 2 -> 3 (w: 2.70) | [Poprawiono ,] 
3        | Krawedz krawedz3: 3 -> 1 (w: 0.90) | [Poprawiono ,] 
===============================================================
INFO: Dokonano 3 autokorekt separatora dziesiętnego.

Wczytywanie zakończone sukcesem. Dane są poprawne technicznie.
Statystyki: 3 węzłów, 3 krawędzi.
---------------------------------------------------------------

\end{lstlisting}
\newpage

\subsubsection{Obsługa pustych linii i białych znaków (\texttt{test\_puste.txt})}
Parser musi ignorować wiersze niezawierające danych, aby nie przerywać procesu wczytywania grafu.
\textbf{Fragment struktury:} \begin{verbatim}
k1;1;2;1.0

k2;2;3;1.5

k3;3;1;0.8
\end{verbatim}
\begin{lstlisting}[language=bash, caption={Log terminala - pomijanie pustych linii}]
Uruchamiam przetwarzanie...

===============================================================
 RAPORT ANALIZY DANYCH WEJSCIOWYCH 
===============================================================
Linia    | Zinterpretowane dane                          | Status         
---------------------------------------------------------------
1        | Krawedz k1: 1 -> 2 (w: 1.00)                  | [OK]           
2       | [INFO] Pominięto pustą linię.
3        | Krawedz k2: 2 -> 3 (w: 1.50)                  | [OK]           
4       | [INFO] Pominięto pustą linię.
5        | Krawedz k3: 3 -> 1 (w: 0.80)                  | [OK]           
===============================================================
Linia 2: [INFO] Napotkano pustą linię - pomijam.
Linia 4: [INFO] Napotkano pustą linię - pomijam.

Wczytywanie zakończone sukcesem. Dane są poprawne technicznie.
Statystyki: 3 węzłów, 3 krawędzi.
---------------------------------------------------------------

\end{lstlisting}
\subsection{Pozostałe zweryfikowane mechanizmy}

Poza testami opartymi na plikach tekstowych, aplikacja została poddana weryfikacji funkcjonalnej w następujących obszarach:

\begin{itemize}
    \item \textbf{Parametry CLI}: Sprawdzono reakcję na brak flag \texttt{-i}, \texttt{-o} oraz wpisanie nieistniejących algorytmów. Program każdorazowo przerywa pracę i wyświetla instrukcję obsługi \texttt{Usage}.
    \item \textbf{Logika struktur danych}: Przetestowano mechanizm odrzucania duplikatów krawędzi oraz ignorowania pętli własnych, co zapewnia stabilność obliczeniową.
    \item \textbf{Formaty wyjściowe}: Zweryfikowano poprawność zapisu binarnego (\texttt{.bin}) oraz tekstowego przy użyciu funkcji \texttt{save\_graph}.
    \item \textbf{Bezpieczeństwo I/O}: Program poprawnie identyfikuje brak pliku wejściowego (Kod 2) oraz brak uprawnień do zapisu.
\end{itemize}

\subsection{Test pozytywny – wzorcowy przebieg algorytmów: Tutte'a oraz Fruchtermanna}

Do weryfikacji pełnej ścieżki krytycznej programu użyto poprawnego, spójnego grafu planarnego (cykl o 3 wierzchołkach). Poniżej przedstawiono wyniki testów dla obu zaimplementowanych metod rozmieszczania.

\textbf{Zawartość pliku wejściowego (\texttt{poprawny\_plik.txt}):}
\begin{verbatim}
K1;1;2;1.0
K2;2;3;1.0
K3;3;1;1.0
\end{verbatim}

\subsubsection{Logi z konsoli (Przebieg operacji)}
Poniższe listingi przedstawiają raporty walidacyjne oraz potwierdzenie sukcesu obliczeń dla obu dostępnych silników graficznych.

\begin{lstlisting}[language=bash, caption={Przebieg programu dla algorytmu Tutte'a (metoda deterministyczna)}]
Uruchamiam przetwarzanie...

===============================================================
 RAPORT ANALIZY DANYCH WEJSCIOWYCH 
===============================================================
Linia    | Zinterpretowane dane                          | Status         
---------------------------------------------------------------
1        | Krawedz k1: 1 -> 2 (w: 1.00)                  | [OK]           
2        | Krawedz k2: 2 -> 3 (w: 1.00)                  | [OK]           
3        | Krawedz k3: 3 -> 1 (w: 1.00)                  | [OK]           
===============================================================

Wczytywanie zakończone sukcesem. Dane są poprawne technicznie.
Statystyki: 3 węzłów, 3 krawędzi.
---------------------------------------------------------------

Sprawdzanie spójności grafu za pomocą algorytmu DFS ...
Sukces: Graf jest spójny. Przechodzę do obliczeń.

Uruchamiam algorytm Tutte...

Sprawdzanie planarności wygenerowanego układu (Tutte)...
Sukces! Graf planarny (Tutte).

---------------------------------------------------------------
>>> SUKCES: Wynik został zapisany do pliku [wynik.txt] <<<
>>> INFO: Nie podano parametru -f bin. Zapisano domyślnie jako plik TEKSTOWY (.txt) <<<
\end{lstlisting}
\newpage
\begin{lstlisting}[language=bash, caption={Przebieg programu dla algorytmu Fruchtermanna-Reingolda (metoda siłowa)}]
Uruchamiam przetwarzanie...

===============================================================
 RAPORT ANALIZY DANYCH WEJSCIOWYCH 
===============================================================
Linia    | Zinterpretowane dane                          | Status         
---------------------------------------------------------------
1        | Krawedz k1: 1 -> 2 (w: 1.00)                  | [OK]           
2        | Krawedz k2: 2 -> 3 (w: 1.00)                  | [OK]           
3        | Krawedz k3: 3 -> 1 (w: 1.00)                  | [OK]           
===============================================================

Wczytywanie zakończone sukcesem. Dane są poprawne technicznie.
Statystyki: 3 węzłów, 3 krawędzi.
---------------------------------------------------------------

Sprawdzanie spójności grafu za pomocą algorytmu DFS ...
Sukces: Graf jest spójny. Przechodzę do obliczeń.

Uruchamiam symulację Fruchtermana (max 3000 prób)...

Sprawdzanie planarności wygenerowanego układu (Fruchterman)...
Sukces! Graf planarny (Fruchterman). Wygenerowany po 1 probach.

---------------------------------------------------------------
>>> SUKCES: Wynik został zapisany do pliku [wynik.txt] <<<
>>> INFO: Nie podano parametru -f bin. Zapisano domyślnie jako plik TEKSTOWY (.txt) <<<
\end{lstlisting}
\newpage
\section{Interpretacja wyników i wnioski}

\subsection{Algorytm Tutte’a (Deterministyczny)}
\begin{itemize}
    \item \textbf{Determinizm obliczeniowy:} Uzyskanie identycznych współrzędnych w każdej próbie (np. ID 1: 833.33, 500.00) potwierdza, że algorytm poprawnie rozwiązuje układ równań liniowych, dążąc do jedynego rozwiązania optymalnego.
    \item \textbf{Geometria układu:} Rozłożenie wierzchołków na planie trójkąta (rama wypukła) jest zgodne z założeniami teoretycznymi dla grafów o małej liczbie węzłów, gwarantując brak przecięć krawędzi wewnątrz ramy.
    \item \textbf{Sprawność modułów:} Poprawna struktura współrzędnych w pliku wynikowym potwierdza bezbłędną komunikację między modułem matematycznym a eksporterem danych.
\end{itemize}

\begin{lstlisting}[language={}, caption={Wynik algorytmu Tutte'a (zawsze powtarzalny)}]
1 833.33 500.00
2 500.00 833.33
3 166.67 500.00
\end{lstlisting}

\subsection{Algorytm Fruchtermana-Reingolda (Stochastyczny)}
\begin{itemize}
    \item \textbf{Niedeterminizm pożądany:} Zmienność wyników (np. X=195.03 vs X=555.27 dla ID 1) potwierdza poprawne działanie mechanizmu stochastycznego i inicjalizacji losowej pozycji wierzchołków.
    \item \textbf{Minimalizacja energii:} Uzyskanie czytelnego układu przy różnych współrzędnych dowodzi, że algorytm skutecznie odnajduje lokalne minima energii układu sił przyciągania i odpychania.
    \item \textbf{Zgodność z modelem:} Proces symulacji fizycznej przebiega zgodnie z założeniami modelu siłowego, oferując elastyczne i estetyczne rozmieszczenie grafu.
\end{itemize}

\begin{lstlisting}[language={}, caption={Wynik Fruchtermana - Proba 1}]
1 195.03 370.56
2 540.33 829.28
3 764.95 300.88
\end{lstlisting}

\begin{lstlisting}[language={}, caption={Wynik Fruchtermana - Proba 2}]
1 555.27 826.83
2 755.76 288.82
3 189.58 384.19
\end{lstlisting}




\subsection{Odporność systemu i walidacja}
\begin{itemize}
    \item \textbf{Wielowarstwowa walidacja:} Program skutecznie filtruje dane na trzech poziomach:
        \begin{itemize}
            \item \textbf{Syntaktycznym:} Blokada znaków ASCII, spacji i błędnych nazw.
            \item \textbf{Semantycznym:} Kontrola zakresów ID (1--1000) oraz wag (0.1--100.0).
            \item \textbf{Topologicznym:} Weryfikacja spójności grafu algorytmem DFS.
        \end{itemize}
    \item \textbf{Stabilność:} Mechanizmy obronne eliminują ryzyko wystąpienia błędów krytycznych (np. \textit{segmentation fault}) przy podaniu uszkodzonych plików.
    \item \textbf{Autokorekta:} Moduł \textit{sanitize} poprawnie normalizuje dane (zamiana przecinków na kropki), zwiększając tolerancję na błędy wprowadzania.
\end{itemize}

\subsection{Porównanie i konkluzje końcowe}
\begin{table}[h]
\centering
\begin{tabular}{|l|p{5cm}|p{5cm}|}
\hline
\textbf{Cecha} & \textbf{Algorytm Tutte’a} & \textbf{Algorytm Fruchtermana} \\ \hline
\textbf{Powtarzalność} & Całkowita (determinizm) & Brak (wynik zależny od losowania) \\ \hline
\textbf{Zaleta} & Gwarancja planarności, szybkość & Naturalny i estetyczny układ \\ \hline
\textbf{Zastosowanie} & Grafy o sztywnej strukturze & Dowolne grafy spójne \\ \hline
\end{tabular}
\end{table}

\textbf{Wniosek końcowy:} Przeprowadzone testy potwierdzają pełną sprawność oprogramowania. Program jest odporny na błędy użytkownika, a zaimplementowane algorytmy realizują swoje cele matematyczne i wizualne zgodnie ze specyfikacją projektową.

\subsection{Wnioski z testów}
Przeprowadzone testy potwierdzają, że program jest odporny na błędy użytkownika oraz nieprawidłowe formatowanie danych. Wielowarstwowa walidacja (od znaków ASCII po spójność topologiczną) gwarantuje, że algorytmy pracują na poprawnych danych.

\end{document}
